\section{The cubic families of type \{6,6\}}\label{the-reused-66-calculation}

We incorporate the earlier cubic proof {[}C66{]} here, including its source-word calculation. It uses

\[a(x,y,z)=(y,x,z),\quad b(x,y,z)=(1-x,z,y),\quad
c(x,y,z)=(-y,-x,z).\tag{17}\]

Its point quotient is the same signed permutation group of order 48;
its translation kernel is \(2D_3\), of index sixteen in \(\mathbb Z^3\).
The retained arithmetic checker records raw affine differences. Its
translation-lattice conclusion uses the invariant-lattice version of
Lemma AL.1: \(2D_3\) is signed-permutation invariant, contains every raw
difference, and is generated by exhibited raw differences. No general
claim that arbitrary raw differences alone generate a kernel is used.
An exact index48 kernel presentation identifies the complete universal
group

\[\langle a,b,c\mid a^2,b^2,c^2,(ab)^6,(bc)^6,(ac)^2,(cbab)^3\rangle
\tag{18}\]

with this affine group. For clarity, the three triples of period words are displayed in full. Put
\[
 z=((bc)^3(ab)^3)^2.
\]
The triples $(u_i,v_i,w_i)$ of the source's Theorem 5.1 are
\[
\begin{array}{c|ccc}
i&u_i&v_i&w_i\\\hline
1&((acb)^4)^b&(cba)^4&((cba)^4)^b\\
2&z^{ab}&z^c&z^{cb}\\
3&(((ab)^2(cb)^2)^3)^a&((bc)^2(ba)^2)^3&(((ba)^2(bc)^2)^3)^c .
\end{array}
\]
Repeated affine multiplication gives the following complete translation table:
\[
\begin{array}{c|ccc|c|c}
i&u_i&v_i&w_i&\Lambda_i&[L:\Lambda_i]\\\hline
1&(4,0,0)&(0,-4,0)&(0,0,-4)&4\mathbb Z^3&4\\
2&(4,-4,0)&(4,0,-4)&(-4,-4,0)&4D_3&8\\
3&(4,-4,4)&(4,-4,-4)&(-4,-4,4)&8\mathbb Z^3+4\mathbb Z(1,1,1)&16 .
\end{array}
\]
Here $D_3=\{(x,y,z)\in\mathbb Z^3:x+y+z\equiv0\pmod2\}$.
For example, conjugating a translation $t_v$ by $(M,u)$ changes its vector to $M^{-1}v$; thus the conjugation exponents in the displayed words are evaluated with the same convention as the source. The integer-only program evaluates all nine words, not merely their determinant indices. Those indices in $\mathbb Z^3$ are $64,128,256$, while $[\mathbb Z^3:L]=16$. Each displayed triple lies in the proposed lattice and has determinant equal to its index, proving that it is a basis. Each lattice is invariant under coordinate permutations and independent sign changes, so is normal in the affine group.

The source's three presentations adjoin $u_i^m=v_i^m=w_i^m=1$ to (18). By Theorem AL.3 their groups are exactly $E/(m\Lambda_i)$, for every positive integer $m$. In particular their orders are $192m^3,384m^3,768m^3$. This is an identification of each entire presentation, with its distinguished marking; an affine quotient of matching numerical type would not suffice.

Both adjacent parabolics have order twelve and faithful linear image.
Their linear intersection has four elements. Two lift to the common
identity and middle involution; the other two differ by translation
\((0,2,2)\). None of the three period lattices, or their positive
multiples, contains that vector, since they lie in \(4\mathbb Z^3\).
This is the complete all-parameter IP check. The quotients are solvable
by the same point-group argument as before. It proves the exact
source Theorem 5.1, with its presented groups, rather than another
family sharing their numerical parameters.

Taking \(m=2^j\) covers all exponents \(n\geq6\) in order \(3\cdot2^n\).
For \(n=5\), take the tetrahedral marking
\((12),(23),(34)\) in \(S_4\), and attach independent central involutions
\(u,v\) to the first and third markers. In the tetrahedral marking the first adjacent subgroup consists of all permutations fixing 4, and the second of all permutations fixing 1. Their intersection is precisely the order-two subgroup on $(23)$, so the rank-three test proves its IP directly. The resulting marking in
\(S_4\times\langle u,v\rangle\) has adjacent products of order six.
Their cubes recover \(u,v\), so it generates all 96 elements. Its two
parabolic kernels are respectively \(\langle u\rangle\) and
\(\langle v\rangle\); the kernel criterion preserves the tetrahedral
intersection. This proves the full source Corollary 5.2. The proof and
its attribution are reused from {[}C66{]}.
